% Replacements for the output-state-space section of the final draft.
% 1. Insert before the compactness/convexity paragraph.
Since $c_t(0)=t$ and $k^2t>1$, the origin does not belong to
$\mathscr D_{k,t}$.  Compactness and nonnegativity of the coordinates
therefore give
\begin{equation}
    m_{k,t}:=\min_{u\in\mathscr D_{k,t}}\sum_i u_i>0.
    \label{eq:positive-total-mass}
\end{equation}

% 2. The Hausdorff-distance definition and support-duality lemma should
%    both explicitly require nonempty compact sets.  Replace the
%    opening sentence of Lemma~\ref{lem:hausdorff-support} by:
For nonempty compact convex sets $C,K\subseteq\mb M_k^{\rm sa}$,
\begin{equation}
    d_H^{(1)}(C,K)
    =\sup_{\|H\|_\infty\le1}\bigl|h_C(H)-h_K(H)\bigr|.
\end{equation}

% 3. Optional shorter replacement for the proof of the output-space theorem.
\begin{proof}[Proof of Theorem~\ref{thm:output-space-limit}]
Write $h_n:=h_{\mathscr C_n}$ and $h:=h_{\mathscr K_{k,t}}$.
For $H=\diag(a)$, von Neumann's trace inequality and permutation
invariance give
\begin{equation}
    h(H)=\max_{u\in\mathscr D_{k,t}}
       \frac{\sum_i a_i u_i}{\sum_i u_i}.
\end{equation}
Put $\lambda_*=h(H)$ and
\begin{equation}
    f_n(\lambda):=\lambda_{\max}(S_n(a-\lambda\mathbf1)),
    \qquad
    f(\lambda):=\max_{u\in\mathscr D_{k,t}}
       \sum_i(a_i-\lambda)u_i.
\end{equation}
By the random compression formula, $f_n(\lambda)\to f(\lambda)$
almost surely, simultaneously in $a$ and $\lambda$.  Since the
maximum defining $\lambda_*$ is attained, every $\varepsilon>0$
satisfies
\begin{equation}
    f(\lambda_*+\varepsilon)\le-\varepsilon m_{k,t}<0
    <\varepsilon m_{k,t}\le f(\lambda_*-\varepsilon).
\end{equation}
Lemma~\ref{lem:support-channel} gives
$h_n(H)\le\lambda\iff f_n(\lambda)\le0$.  The strict signs therefore
imply $\lambda_*-\varepsilon<h_n(H)\le\lambda_*+\varepsilon$ for all
large $n$, so $h_n(H)\to h(H)$.

For fixed $H=V\diag(a)V^*$, apply the same argument to
$(I_n\otimes V^*)P_n(I_n\otimes V)$, which has the same distribution
as $P_n$.  Thus $h_n(H)\to h(H)$ almost surely for each fixed $H$.
Intersect these events over a countable dense subset of the
operator-norm unit ball.  Both $h_n$ and $h$ are $1$-Lipschitz there,
so a finite $\delta$-net bounds the uniform error by the largest
error on the net plus $2\delta$.  Compactness then gives uniform
convergence.  Lemma~\ref{lem:hausdorff-support} completes the proof.
\end{proof}
